% Kreis – bank/kreis/ (zusammenbau v0.4, Heft)
\einheitenkopf[t6]{Kreis}
\setcounter{aufgabe}{6}

\begin{aufgabe}{Kreisteile – Prüfungsaufgaben}
\begin{teile}
\teil Im Bild ist ein Kreisausschnitt grau gefärbt. Sein Mittelpunktswinkel beträgt $130^\circ$, der Winkel des weißen Teils heißt $\alpha$. Wie viel Prozent der Kreisfläche sind grau? Runde auf eine Stelle \hfill (P10 2025 OS).

\kreissektor{130}{$130^\circ$}
≈ \leerfeld[\%]
\teil Ein rundes Beet ist geteilt: Auf dem grauen Ausschnitt im Bild mit dem Mittelpunktswinkel $155^\circ$ wachsen Rosen, auf dem weißen Teil mit dem Winkel $\alpha$ Gras. Wie viel Prozent des Beets tragen Rosen? Runde auf eine Stelle \hfill (P10 2025 OS).

\kreissektor{155}{$155^\circ$}
≈ \leerfeld[\%]
\end{teile}
\end{aufgabe}
